% Unit 079 terjemahan Bahasa Indonesia dari chapter6.tex baris 1361--1539.
% Sumber: Methods of Algebra, Volume 2, commit
% 9a5803ff77dd3257484cb177f851a73770a59dd3 (CC BY 4.0).
% stable-unit-id: o014.aljabr2.chapter6.finite-index-subgroups
% Cakupan: Bagian 6.8 lengkap, ``Subgrup Berindeks Hingga'';
% berhenti tepat sebelum Bagian 6.9 pada baris 1540.

% segment-id: o014.li2.chapter6.finite-index-subgroups.h001
\section{Subgrup Berindeks Hingga}\label{sec:group-finite-index}
\hypertarget{o014.aljabr2.chapter6.finite-index-subgroups}{ }
% segment-id: o014.li2.chapter6.finite-index-subgroups.p001
Dalam bagian ini, dengan gelanggang komutatif $\Bbbk$ tetap, kita meninjau grup
$G$ beserta subgrupnya $H$. Tujuannya ialah membahas beberapa fenomena khusus
ketika indeks $(G:H)$ hingga.

% segment-id: o014.li2.chapter6.finite-index-subgroups.p002
Definisi~\ref{def:induced-module} mendefinisikan modul-$G$ terinduksi
$\Ind^G_H(N)$ dan $\iInd^G_H(N)$ untuk setiap modul-$H$ $N$. Langkah pertama
dalam bagian ini ialah merealisasikan keduanya sebagai ruang pemetaan
tertentu; pada tahap ini, $(G:H)$ belum perlu hingga.

% segment-id: o014.li2.chapter6.finite-index-subgroups.p003
Untuk setiap modul-$\Bbbk$ $V$, berikan struktur modul-$G$ berikut pada
himpunan pemetaan $\{f:G\to V\}$. Struktur modul-$\Bbbk$ berasal dari operasi
titik demi titik pada pemetaan, sedangkan aksi $G$ diberikan oleh
% segment-id: o014.li2.chapter6.finite-index-subgroups.q001
\[ (g f)(x) = f(xg), \quad f: G \to V, \; g, x \in G. \]

% segment-id: o014.li2.chapter6.finite-index-subgroups.p004
Untuk setiap pemetaan $f:G\to V$, tuliskan
$\Supp(f):=\{x\in G:f(x)\neq0\}$. Demi memudahkan pernyataan, pilih wakil
$(g_i)_{i\in I}$ bagi dekomposisi koset sedemikian sehingga
\footnote{Koreksi penerjemah: sumber menuliskan indeks $i\in H$, sedangkan
dekomposisi berikutnya diindeks oleh himpunan $I$.}
% segment-id: o014.li2.chapter6.finite-index-subgroups.q002
\[ G = \bigsqcup_{i \in I} g_i H = \bigsqcup_{i \in I} Hg_i^{-1}. \]

% segment-id: o014.li2.chapter6.finite-index-subgroups.le001
\begin{lema}[Modul terinduksi sebagai ruang pemetaan]\label{prop:Ind-as-mappings}
	Terdapat isomorfisme kanonik modul-$G$ berikut.
% segment-id: o014.li2.chapter6.finite-index-subgroups.q003
	\begin{equation*}\begin{gathered}
% segment-id: o014.li2.chapter6.finite-index-subgroups.g001
		\begin{tikzcd}[column sep=small, row sep=tiny, ampersand replacement=\&]
			\Ind^G_H(N) \arrow[leftrightarrow, r, "\sim"] \& \left\{\begin{array}{r|l}
				f: G \to N & \forall h \in H, \forall x \in G, \\
				& f(hx) = h \cdot f(x)
			\end{array}\right\} \\
			\varphi \arrow[mapsto, r] \& {\varphi|_{G \subset \Bbbk[G]}} \\
			{\left[\sum_g a_g g \mapsto \sum_g a_g f(g)\right]} \& f, \arrow[mapsto, l]
		\end{tikzcd} \\
% segment-id: o014.li2.chapter6.finite-index-subgroups.g002
		\begin{tikzcd}[column sep=small, row sep=tiny, ampersand replacement=\&]
			\iInd^G_H(N) \arrow[leftrightarrow, r, "\sim"] \& \left\{\begin{array}{r|l}
				f: G \to N & \text{seperti di atas, tetapi}\; H \backslash \Supp(f) \;\text{hingga}
			\end{array}\right\} \\
			g_i \otimes y \arrow[mapsto, r] \& {\left[ h g_i^{-1} \mapsto h y, \;\text{bernilai $0$ di luar $Hg_i^{-1}$} \right]} \\
			\sum_i g_i \otimes f(g_i^{-1}) \& f, \arrow[mapsto, l]
		\end{tikzcd}
	\end{gathered}\end{equation*}
	dengan $g\in G$, $y\in N$, dan $i\in I$. Ruas kanan isomorfisme diberi
	struktur modul-$G$ seperti dalam pembahasan sebelumnya. Pemetaan-pemetaan
	ini tidak bergantung pada pilihan wakil $g_i$.
\end{lema}
% segment-id: o014.li2.chapter6.finite-index-subgroups.pf001
\begin{bukti}
	Verifikasilah secara langsung bahwa semua pemetaan terdefinisi dengan baik,
	$G$-ekuivarian, saling invers, dan tidak bergantung pada pilihan wakil.
\end{bukti}

% segment-id: o014.li2.chapter6.finite-index-subgroups.p005
Isomorfisme dalam Proposisi~\ref{prop:pullback-two-adjoint} mudah dituliskan
ulang melalui ruang pemetaan di atas. Sebagai contoh,
% segment-id: o014.li2.chapter6.finite-index-subgroups.g003
\begin{equation}\label{eqn:Ind-as-mappings}
	\begin{tikzcd}[row sep=tiny]
		\Hom_H\left(\Res^G_H M, N \right) \arrow[r, "\sim"] & \Hom_G\left(M, \Ind^G_H N \right) \\
		\theta \arrow[mapsto, r] & {\left[ x \mapsto [g \mapsto \theta(gx)] \right]}, \\
		\Hom_H\left(N, \Res^G_H M\right) \arrow[r, "\sim"] & \Hom_G\left(\iInd^G_H N, M\right) \\
		\psi \arrow[mapsto, r] & {\left[ f \mapsto \sum_{\text{koset}\; gH} g \psi(f(g^{-1})) \right]}.
	\end{tikzcd}
\end{equation}

% segment-id: o014.li2.chapter6.finite-index-subgroups.co001
\begin{korolari}\label{prop:iInd-Ind}
	Untuk setiap modul-$H$ $N$, $\iInd^G_H(N)$ dapat dibenamkan secara kanonik
	sebagai submodul-$G$ dari $\Ind^G_H(N)$. Jika $(G:H)$ hingga, keduanya sama.
\end{korolari}
% segment-id: o014.li2.chapter6.finite-index-subgroups.pf002
\begin{bukti}
	Modul-$G$ pada ruas kanan Lema~\ref{prop:Ind-as-mappings} mempunyai inklusi
	alami. Jika $(G:H)$ hingga, setiap pemetaan $f:G\to N$ memenuhi bahwa
	$H\backslash\Supp(f)$ hingga.
\end{bukti}

% segment-id: o014.li2.chapter6.finite-index-subgroups.p006
Selanjutnya kita mendefinisikan pemetaan korestriksi dengan asumsi bahwa
$(G:H)$ hingga. Untuk menyederhanakan notasi, bagi modul-$G$ $M$, modul-$H$
$\Res^G_H(M)$ berikut ini cukup dituliskan sebagai $M$.

% segment-id: o014.li2.chapter6.finite-index-subgroups.l001
\begin{itemize}
% segment-id: o014.li2.chapter6.finite-index-subgroups.i001
	\item Keluarga homomorfisme kanonik dalam
	\eqref{eqn:change-of-group-coh}, dengan $\varphi$ diambil sebagai inklusi
	$H\hookrightarrow G$, disebut \emph{restriksi} dan dituliskan
% segment-id: o014.li2.chapter6.finite-index-subgroups.q004
	\[ \mathrm{res}^n: \Hm^n(G, M) \to \Hm^n(H, M), \quad n \in \Z_{\geq 0}. \]
% segment-id: o014.li2.chapter6.finite-index-subgroups.i002
	\item Demikian pula, dalam kasus homologi terdapat homomorfisme kanonik
	dengan arah sebaliknya,
% segment-id: o014.li2.chapter6.finite-index-subgroups.q005
	\[ \mathrm{res}_n: \Hm_n(H, M) \to \Hm_n(G, M), \quad n \in \Z_{\geq 0}. \]
\end{itemize}
% segment-id: o014.li2.chapter6.finite-index-subgroups.p007
Jika $\mathrm{res}^n$ dideskripsikan melalui kompleks kokrantai standar
dengan Proposisi~\ref{prop:grp-equiv}, efeknya hanyalah membatasi pemetaan
$G^n\to M$ ke $H^n$; dari sinilah namanya berasal. Bagian
\S\ref{sec:change-of-group} telah mengangkat operasi-operasi ini ke
kategori turunan.
\index[sym1]{res-n@$\mathrm{res}_n$, $\mathrm{res}^n$}

% segment-id: o014.li2.chapter6.finite-index-subgroups.d001
\begin{definisi}\label{def:cor-zero}
	\index[sym1]{nu-GH@$\nu_{G\mid H}$, $\nu^{G\mid H}$}
	Misalkan $(G:H)$ hingga. Untuk setiap modul-$G$ $M$, definisikan
	homomorfisme modul-$\Bbbk$ berikut:
% segment-id: o014.li2.chapter6.finite-index-subgroups.g004
	\[\begin{tikzcd}[row sep=tiny]
		\nu^{G|H}: M^H \arrow[r] & M^G \\
		x \arrow[mapsto, r] & \sum_{\bar{g} \in G/H} gx, \\
		\nu_{G|H}: M_G \arrow[r] & M_H \\
		(x \;\text{citranya}) \arrow[mapsto, r] & \sum_{\bar{g} \in H \backslash G} (gx \;\text{citranya}),
	\end{tikzcd}\]
	dengan $g\in G$ sembarang wakil dari koset $\bar g$. Kedua homomorfisme
	bersifat funktorial dalam $M$.
\end{definisi}

% segment-id: o014.li2.chapter6.finite-index-subgroups.p008
Dalam baris pertama definisi, $gx$ jelas hanya bergantung pada koset
$\bar g=gH$. Baris kedua memerlukan sedikit penjelasan. Untuk $x\in M$ yang
diberikan, citra $gx$ dalam $M_H$ hanya bergantung pada koset $Hg$. Karena
itu, untuk setiap koset $t\in H\backslash G$, kita dapat memilih sembarang
wakil $\theta(t)\in G$ dan menjumlahkan citra-citra $\theta(t)x$, dengan $t$
merentang $H\backslash G$. Jika $x$ diganti dengan $g'x$, dengan
$g'\in G$, maka $(\theta(t)g')_{t\in H\backslash G}$ tetap merupakan
keluarga wakil yang menjelajahi $H\backslash G$, sehingga
% segment-id: o014.li2.chapter6.finite-index-subgroups.q006
\[ \sum_t (\theta(t) g' x \;\text{citranya}) = \sum_t (\theta(t) x \;\text{citranya}). \]
Jadi $\nu_{G|H}$ terdefinisi dengan baik pada tingkat $M_G$.

% segment-id: o014.li2.chapter6.finite-index-subgroups.dp001
\begin{definisi-proposisi}\label{def:cor}
	\index[sym1]{cor-n@$\mathrm{cor}_n$, $\mathrm{cor}^n$}
	\index{korestriksi (corestriction)}
	Misalkan $(G:H)$ hingga dan $M$ modul-$G$.
% segment-id: o014.li2.chapter6.finite-index-subgroups.l002
	\begin{enumerate}[(i)]
% segment-id: o014.li2.chapter6.finite-index-subgroups.i003
		\item Terdapat keluarga homomorfisme kanonik
% segment-id: o014.li2.chapter6.finite-index-subgroups.q007
		\[ \mathrm{cor}^n: \Hm^n(H, M) \to \Hm^n(G, M), \quad n \in \Z_{\geq 0}, \]
		yang kompatibel dengan barisan eksak panjang dan pada derajat $0$
		memberikan $\nu^{G|H}:M^H\to M^G$.
% segment-id: o014.li2.chapter6.finite-index-subgroups.i004
		\item Demikian pula, terdapat keluarga homomorfisme kanonik yang
		kompatibel dengan barisan eksak panjang,
% segment-id: o014.li2.chapter6.finite-index-subgroups.q008
		\[ \mathrm{cor}_n: \Hm_n(G, M) \to \Hm_n(H, M), \]
		yang pada derajat $0$ memberikan $\nu_{G|H}:M_G\to M_H$.
	\end{enumerate}
	Kedua keluarga homomorfisme itu disebut \emph{korestriksi}, karena arahnya
	berlawanan dengan restriksi\footnote{Dengan alasan yang cukup kuat, banyak
	literatur menyebut homomorfisme korestriksi sebagai ``transfer''. Buku ini
	memakai istilah homomorfisme transfer dalam arti yang lebih sempit; lihat
	pembahasan selanjutnya.}.
\end{definisi-proposisi}
% segment-id: o014.li2.chapter6.finite-index-subgroups.pf003
\begin{bukti}
	Pertama, tinjau kasus kohomologi. Kita berikan dua penjelasan. Pertama,
	$\Res^G_H:G\dcate{Mod}\to H\dcate{Mod}$ merupakan funktor eksak yang
	mempertahankan objek injektif. Karena itu, $M\mapsto\Hm^n(H,M)$ merupakan
	funktor turunan kanan ke-$n$ dari $M\mapsto M^H$, dengan
	$n\in\Z_{\geq0}$. Funktorialitas $\nu^{G|H}$ membuatnya menginduksi
	morfisme funktor-$\delta$ kohomologis
	$\Hm^n(H,\cdot)\to\Hm^n(G,\cdot)$.
	
	Penjelasan kedua melibatkan kategori turunan. Argumen di atas mudah
	diangkat ke tingkat kategori turunan, tetapi di sini kita berikan sudut
	pandang lain berdasarkan adjungsi. Pertama, kita nyatakan bahwa adjungsi
	antara $\iInd^G_H$ dan $\Res^G_H$ terangkat ke tingkat $\RHom$, sehingga
	memberikan
% segment-id: o014.li2.chapter6.finite-index-subgroups.q009
	\[ \RHom_G\left(\iInd^G_H(\Bbbk), M\right) \simeq \RHom_H(\Bbbk, M). \]
	Gagasan serupa telah muncul dalam bukti Teorema~\ref{prop:Shapiro}.
	Kuncinya ialah menafsirkan $\iInd^G_H$ sebagai funktor perubahan gelanggang
	$\Bbbk[H]\dcate{Mod}\to\Bbbk[G]\dcate{Mod}$ yang diinduksi oleh
	$\Bbbk[H]\to\Bbbk[G]$, memperhatikan bahwa funktor tersebut eksak
	(Proposisi~\ref{prop:Ind-preservation}), lalu menerapkan
	Contoh~\ref{eg:change-of-ring-adjunction}.
	
	Selanjutnya, syarat bahwa $(G:H)$ hingga, bersama
	Korolari~\ref{prop:iInd-Ind}, memberikan
% segment-id: o014.li2.chapter6.finite-index-subgroups.q010
	\[ \RHom_G\left(\iInd^G_H(\Bbbk), M\right) \simeq \RHom_G\left( \Ind^G_H(\Bbbk), M \right) \to \RHom_G(\Bbbk, M), \]
	dengan bagian terakhir berasal dari $\Bbbk\to\Ind^G_H(\Bbbk)$. Ini
	merupakan unit pasangan adjoin $(\Res^G_H,\Ind^G_H)$ pada $\Bbbk$, dan
	dapat pula dideskripsikan lebih konkret melalui
	Lema~\ref{prop:Ind-as-mappings}: unsur $t\in\Bbbk$ dipetakan ke pemetaan
	konstan $G\to\Bbbk$ yang bersesuaian.
	
	Dengan demikian diperoleh morfisme kanonik
	$\RHom_H(\Bbbk,M)\to\RHom_G(\Bbbk,M)$. Mengambil $\Hm^n$ pada kedua ruas
	memberikan hasil yang dicari. Menentukan perilakunya untuk $n=0$ sama dengan
	mengganti $\RHom$ oleh $\Hom$ dalam operasi di atas. Secara terperinci,
	jika $\iInd^G_H\subset\Ind^G_H$ direalisasikan sebagai ruang pemetaan dan
	isomorfisme adjungsi dideskripsikan dengan
	\eqref{eqn:Ind-as-mappings}, diperoleh
% segment-id: o014.li2.chapter6.finite-index-subgroups.g005
	\[\begin{tikzcd}[row sep=tiny, column sep=small]
		\Hom_H(\Bbbk, M) \arrow[r, "\sim"] & \Hom_G\left(\iInd^G_H \Bbbk, M\right) \arrow[equal, r] & \Hom_G\left(\Ind^G_H \Bbbk, M\right) \arrow[r] & \Hom_G(\Bbbk, M) \\
		{[1 \mapsto x]} \arrow[mapsto, rr] & & {[f \mapsto \sum_{gH} f(g^{-1}) gx]} \arrow[mapsto, r] & {[1 \mapsto \sum_{gH} gx]}.
	\end{tikzcd}\]
	\footnote{Koreksi penerjemah: sumber menuliskan $\Res^G_M$ dan
	$\Hom_H(\Ind^G_H\Bbbk,M)$; adjungsi dan jenis modul menuntut
	$\Res^G_H$ serta $\Hom_G(\Ind^G_H\Bbbk,M)$.}
	Komposisi di atas jelas sama dengan $\nu^{G|H}:M^H\to M^G$. Terbukti.
\end{bukti}

% segment-id: o014.li2.chapter6.finite-index-subgroups.pr001
\begin{proposisi}\label{prop:res-cor}
	Misalkan $(G:H)$ hingga. Untuk setiap modul-$G$ $M$ dan
	$n\in\Z_{\geq0}$, komposisi homomorfisme berikut,
% segment-id: o014.li2.chapter6.finite-index-subgroups.q011
	\begin{gather*}
		\Hm^n(G, M) \xrightarrow{\mathrm{res}^n} \Hm^n(H, M) \xrightarrow{\mathrm{cor}^n} \Hm^n(G, M), \\
		\Hm_n(G, M) \xrightarrow{\mathrm{cor}_n} \Hm_n(H, M) \xrightarrow{\mathrm{res}_n} \Hm_n(G, M)
	\end{gather*}
	masing-masing merupakan endomorfisme perkalian dengan $(G:H)$.
\end{proposisi}
% segment-id: o014.li2.chapter6.finite-index-subgroups.pf004
\begin{bukti}
	Pertama, tinjau kasus $n=0$. Dalam kasus kohomologi, pemetaannya ialah
	$M^G\hookrightarrow M^H\xrightarrow{\nu^{G|H}}M^G$, yang jelas memetakan
	$x\in M^G$ ke $(G:H)x$. Dalam kasus homologi, pemetaannya ialah
	$M_G\xrightarrow{\nu_{G|H}}M_H\twoheadrightarrow M_G$; pemeriksaan langsung
	memberikan kesimpulan yang sama.
	
	Sekarang misalkan $n\geq1$. Dalam kasus kohomologi, ambil pembenaman
	$M\hookrightarrow I$ dengan $I$ modul-$G$ injektif; maka $I$ juga modul-$H$
	injektif. Tuliskan $L:=I/M$. Funktorialitas restriksi dan korestriksi,
	bersama barisan eksak panjang, memberikan diagram komutatif dengan
	kolom-kolom eksak
% segment-id: o014.li2.chapter6.finite-index-subgroups.g006
	\[\begin{tikzcd}
		\Hm^{n-1}(G, L) \arrow[r, "{\mathrm{res}^{n-1}}"] \arrow[twoheadrightarrow, d] & \Hm^{n-1}(H, L) \arrow[r, "{\mathrm{cor}^{n-1}}"] \arrow[twoheadrightarrow, d] & \Hm^{n-1}(G, L) \arrow[twoheadrightarrow, d] \\
		\Hm^n(G, M) \arrow[r, "{\mathrm{res}^n}"] & \Hm^n(H, M) \arrow[r, "{\mathrm{cor}^n}"] & \Hm^n(G, M),
	\end{tikzcd}\]
	Secara rekursif, komposisi baris pertama merupakan perkalian dengan $(G:H)$.
	Karena semua panah vertikal surjektif, demikian pula komposisi baris kedua.
	Argumen untuk homologi serupa: cukup ambil epimorfisme $P\twoheadrightarrow M$
	dengan $P$ modul-$G$ projektif.
	
	Perhatikan bahwa kasus $n=0$ juga memungkinkan kita menyimpulkan pada
	tingkat kategori turunan bahwa komposisi
	$\RHom_G(\Bbbk,M)\to\RHom_H(\Bbbk,M)\to\RHom_G(\Bbbk,M)$ adalah
	$(G:H)\identity$. Kasus homologi serupa.
\end{bukti}

% segment-id: o014.li2.chapter6.finite-index-subgroups.co002
\begin{korolari}\label{prop:group-coh-torsion}
	Misalkan $|G|=m\in\Z_{\geq1}$. Untuk setiap modul-$G$ $M$, berlaku
% segment-id: o014.li2.chapter6.finite-index-subgroups.q012
	\[ n \geq 1 \implies m \Hm^n(G, M) = m \Hm_n(G, M) = \{0\}. \]
\end{korolari}
% segment-id: o014.li2.chapter6.finite-index-subgroups.pf005
\begin{bukti}
	Dalam Proposisi~\ref{prop:res-cor}, ambil $H=\{1\}$ dan gunakan fakta
	sederhana berikut: jika $n\geq1$, maka $\Hm^n(\{1\},M)$ dan
	$\Hm_n(\{1\},M)$ keduanya nol.
\end{bukti}

% segment-id: o014.li2.chapter6.finite-index-subgroups.p009
Berikut ini, ambil $\Bbbk=\Z$ dan pandang $\Z$ sebagai modul-$G$ trivial.
Contoh~\ref{eg:group-H1-further} memberikan isomorfisme kanonik
$\Hm_1(G,\Z)\simeq G_{\mathrm{ab}}$ dan
$\Hm_1(H,\Z)\simeq H_{\mathrm{ab}}$. Jika $(G:H)$ hingga, maka
$\mathrm{cor}_1:\Hm_1(G,\Z)\to\Hm_1(H,\Z)$ bersesuaian dengan
homomorfisme grup
% segment-id: o014.li2.chapter6.finite-index-subgroups.q013
\begin{equation*}
	\mathrm{Ver}_{G|H}: G_{\mathrm{ab}} \to H_{\mathrm{ab}},
\end{equation*}
yang disebut \emph{homomorfisme transfer} dari $G$ ke $H$ (bahasa Jerman:
\textit{die Verlagerung}).
\index{transfer (transfer / Verlagerung)}

% segment-id: o014.li2.chapter6.finite-index-subgroups.p010
Intinya ialah memberikan deskripsi teori grup bagi
$\mathrm{Ver}_{G|H}$. Ingat bahwa $G$ beraksi dari kanan pada
$H\backslash G$ melalui perkalian kanan.

% segment-id: o014.li2.chapter6.finite-index-subgroups.pr002
\begin{proposisi}
	Misalkan $(G:H)$ hingga. Pilih wakil dalam $G$ bagi setiap koset dalam
	$H\backslash G$, yang dinyatakan sebagai pemetaan
	$\theta:H\backslash G\to G$. Untuk setiap $s\in G$ dan
	$t\in H\backslash G$, terdapat tepat satu $h_{t,s}\in H$ sedemikian
	rupa sehingga
% segment-id: o014.li2.chapter6.finite-index-subgroups.q014
	\[ \theta(t)s = h_{t, s} \theta(ts). \]
	Homomorfisme transfer yang didefinisikan sebelumnya dapat dideskripsikan
	sebagai
% segment-id: o014.li2.chapter6.finite-index-subgroups.q015
	\[ \mathrm{Ver}_{G|H}(\underbracket{s\;\text{citranya}}_{\in G_{\mathrm{ab}}}) = \prod_{t \in H \backslash G} \underbracket{h_{t, s}\;\text{citranya}}_{\in H_{\mathrm{ab}}}. \]
\end{proposisi}
% segment-id: o014.li2.chapter6.finite-index-subgroups.pf006
\begin{bukti}
	Ambil barisan eksak pendek modul-$G$
	$0\to\mathfrak{I}_G\to\Z[G]\to\Z\to0$, dengan $\mathfrak{I}_G$ ideal
	augmentasi $\Z[G]$ (Contoh~\ref{eg:group-H1}). Inklusi
	$\Z[H]\subset\Z[G]$ memberikan submodul-$H$
	$\mathfrak{I}_H\subset\mathfrak{I}_G$. Karena $\mathrm{cor}_n$
	kompatibel dengan barisan eksak panjang, dan $\Z[G]$ projektif baik sebagai
	modul-$G$ maupun modul-$H$, diperoleh diagram komutatif berikut:
% segment-id: o014.li2.chapter6.finite-index-subgroups.g007
	\[\begin{tikzcd}
		G_{\mathrm{ab}} \arrow[r, "\sim"'] \arrow[dd, "{\mathrm{Ver}_{G|H}}"'] &
		\Hm_1(G, \Z) \arrow[d, "{\mathrm{cor}_1}"'] \arrow[r, "\sim"] & \Hm_0(G, \mathfrak{I}_G) \arrow[d, "{\mathrm{cor}_0}"] \arrow[equal, r] &  \mathfrak{I}_G/\mathfrak{I}_G^2 \arrow[d, "{\nu_{G|H}}"] \\
		& \Hm_1(H, \Z) \arrow[equal, d] \arrow[hookrightarrow, r] & \Hm_0(H, \mathfrak{I}_G) \arrow[equal, r] & \mathfrak{I}_G / \mathfrak{I}_H \mathfrak{I}_G \\
		H_{\mathrm{ab}} \arrow[r, "\sim"] & \Hm_1(H, \Z) \arrow[r, "\sim"] & \Hm_0(H, \mathfrak{I}_H) \arrow[equal, r] \arrow[u] & \mathfrak{I}_H / \mathfrak{I}_H^2 \arrow[u]
	\end{tikzcd}\]

	Kita akan memeriksa kesamaan yang diperlukan dalam
	$\mathfrak{I}_G/\mathfrak{I}_H\mathfrak{I}_G$. Bagian akhir
	Contoh~\ref{eg:group-H1-further} menunjukkan bahwa citra $s\in G$ dalam
	$G_{\mathrm{ab}}$, melalui baris pertama, bersesuaian dengan citra
	$1_G-s$ dalam $\mathfrak{I}_G/\mathfrak{I}_G^2$. Menerapkan
	$\nu_{G|H}$ padanya memberikan
% segment-id: o014.li2.chapter6.finite-index-subgroups.q016
	\begin{align*}
		\sum_{t \in H \backslash G} \theta(t)(1_G - s) & = \sum_t (\theta(t) - \theta(t)s) = \sum_t (\theta(t) - h_{t, s}\theta(ts)) \\
		& = \sum_t \theta(t) - \sum_t h_{t, s} \theta(ts) \\
		& = \sum_t \theta(ts) - \sum_t h_{t, s} \theta(ts) \;\in \mathfrak{I}_G
	\end{align*}
	dalam $\mathfrak{I}_G/\mathfrak{I}_H\mathfrak{I}_G$. Kesamaan terakhir
	berlaku karena, ketika $t$ bervariasi, baik $t$ maupun $ts$ merentang
	$H\backslash G$. Ekspresi terakhir juga sama dengan
	$\sum_t(1_G-h_{t,s})\theta(ts)$.
	
	Terakhir, perhatikan bahwa
	$(1_G-h_{t,s})\theta(ts)\equiv1_G-h_{t,s}
	\pmod{\mathfrak{I}_H\mathfrak{I}_G}$, sedangkan ruas kanan bersesuaian dengan
	citra $h_{t,s}\;\bmod\mathfrak{I}_H$ dalam
	$\mathfrak{I}_H/\mathfrak{I}_H^2$ melalui baris ketiga. Terbukti.
\end{bukti}

% segment-id: o014.li2.chapter6.finite-index-subgroups.p011
Pemetaan yang dideskripsikan di atas,
$s\bmod G_{\mathrm{der}}\mapsto\prod_t h_{t,s}\bmod H_{\mathrm{der}}$,
merupakan konstruksi klasik dalam teori grup. Karena kita telah
mengidentifikasikannya dengan $\mathrm{cor}_1$, kita juga telah membuktikan
bahwa pemetaan ini tidak bergantung pada pilihan $\theta$.
